%!TEX program = xelatex

\documentclass[compress]{beamer}
%--------------------------------------------------------------------------
% Common packages
%--------------------------------------------------------------------------
\usepackage[english]{babel}
\usepackage{pgfpages} % required for notes on second screen
\usepackage{graphicx}
\usepackage{subfigure}
\usepackage{multicol}
\usepackage{multirow}
\def\block(#1,#2)#3{\multicolumn{#2}{c}{\multirow{#1}{*}{$ #3 $}}}

\usepackage{fontspec}

\usepackage{tabularx,ragged2e}
\usepackage{booktabs}

\usepackage{setspace}

\usepackage{gitdags}
\usepackage[normalem]{ulem} % strikeout
%--------------------------------------------------------------------------
% Load theme
%--------------------------------------------------------------------------
\usetheme{hri}

\usepackage{dtklogos} % must be loaded after theme
\usepackage{tikz}
\usetikzlibrary{arrows,shapes,snakes,calc,mindmap,backgrounds,positioning,svg.path}

\graphicspath{{figs/}}

\usepackage[cache]{minted}
\renewcommand{\theFancyVerbLine}{
  \sffamily\textcolor[rgb]{0.5,0.5,0.5}{\scriptsize\arabic{FancyVerbLine}}}

\newminted{java}{frame=lines,
                 linenos=true,
                 fontsize=\small,
                 xleftmargin=1.8em}

\newmintinline[java]{java}{fontsize=\small}

\newminted{sh}{frame=lines,
                 linenos=false,
                 fontsize=\small,
                 xleftmargin=1.8em}

\newmintinline[sh]{sh}{fontsize=\small}

\setbeamercolor{highlightCol}{bg=hriSec3,fg=white}
\newcommand{\highlight}[1]{%
    \vspace{1em}%
    \begin{beamercolorbox}[wd=\linewidth,ht=2ex,dp=0.7ex]{highlightCol}%
    \centering #1%
    \end{beamercolorbox}%
    \vspace{1em}%
}%

\tikzset{temporal/.code args={<#1>#2#3#4}{%
          \temporal<#1>{\pgfkeysalso{#2}}{\pgfkeysalso{#3}}{\pgfkeysalso{#4}} % \pgfkeysalso doesn't change the path
}}


%--------------------------------------------------------------------------
% General presentation settings
%--------------------------------------------------------------------------
\title{\Medium git}
\subtitle{the basics}
\date{12 Feb. 2016}
\author{Séverin Lemaignan}
\institute{Centre for Robotics \& Neural
Systems\\{\Medium Plymouth University}}

%--------------------------------------------------------------------------
% Notes settings
%--------------------------------------------------------------------------
%\setbeameroption{show notes on second screen}
%\setbeameroption{hide notes}

\begin{document}

\maketitle

%%%%%%%%%%%%%%%%%%%%%%%%%%%%%%%%%%%%%%%%%%%%%%%%%%%%%%%%%%%%%%%%%%%%%%%%%%%%%%%

\imageframe{zipfiles.pdf}
\imageframe{dropbox_zipfiles.pdf}

{
    \fullbackground{../figs/facebook-wall.png}
    \begin{frame}[plain]{}
        \only<2>{
            \begin{center}
                \includegraphics[width=\linewidth]{devil}
            \end{center}
        }
    \end{frame}
}

\section{Code versioning}

\begin{frame}{Why code versioning?}

    \begin{itemize}
        \item The history of your development
        \item Compare the current code with an older version
        \item Roll-back to previous version
        \item Experiment without loosing anything
        \item Trace who did what (at the level of the line of code)
        \item Keep things organized
    \end{itemize}
\end{frame}

\begin{frame}{Atomic commits}

    The single most important concept:

    \highlight{Atomic commit}

    \only<1>{
        A (typically small) commit that represent a {\Medium single, coherent \&
    complete} functional change.
    }
    \only<2>{
    \begin{itemize}
        \item Easy to understand the change
        \item Debugging made easy (\texttt{git bisect})
        \item Collaboration made easy (less, smaller conflict)
        \item Easy to write a useful commit message
    \end{itemize}
    }

\end{frame}

\begin{frame}[fragile]{}

\centering

\begin{tikzpicture}[
    >=latex,
    every edge/.style={draw,thick,hriSec3},
    file/.style={align=left, inner sep=1pt,font=\scriptsize\tt}
]

\node at (-2, 1.5)[file,temporal=<4>{red}{black}{white}] (file1) {\sout{main.cpp}};
\node [file,temporal=<4>{red}{black}{white},below=0.0 of file1.south east,anchor=north east] (file2) {src/main.cpp};
\node [file,temporal=<2>{red}{black}{white},below=0.0 of file2.south east, anchor=north east] (file3) {src/position.cpp};
\node [file,temporal=<2>{red}{black}{white},below=0.0 of file3.south east,anchor=north east] (file4) {src/position.hpp};
\node [file,temporal=<3>{red}{black}{white},below=0.0 of file4.south east,anchor=north east] (file5) {share/model.csv};


    \onslide<2>{
        \gitDAG[grow right sep = 2em]{
            A;
        };

        \gitbranch
        {master}     % node name and text 
        {above=of A} % node placement
        {A}          % target

        \gitHEAD
        {above=of master} % node placement
        {master}          % target


    \path[draw] (file3.east) edge [->, bend left] (A);
    \path[draw] (file4.east) edge [->, bend left] (A);

    }
    \onslide<3>{
        \gitDAG[grow right sep = 2em]{
            A -- B;
        };

        \gitbranch
        {master}     % node name and text 
        {above=of B} % node placement
        {B}          % target

        \gitHEAD
        {above=of master} % node placement
        {master}          % target

        \path[draw] (file5.east) edge [->, bend left,hriSec2] (B);
    }

    \onslide<4->{
        \gitDAG[grow right sep = 2em]{
            A -- B -- C;
        };

        \gitbranch
        {master}     % node name and text 
        {above=of C} % node placement
        {C}          % target

        \gitHEAD
        {above=of master} % node placement
        {master}          % target

        \path[draw] (file1.east) edge [->, bend left,hriSec1Comp] (C);
        \path[draw] (file2.east) edge [->, bend left,hriSec1Comp] (C);
    }

\end{tikzpicture}

\vspace{3em}
\small
\only<2>{
\texttt{git add src/position.*}\\
\texttt{git commit -m"Fix computation of position (float->double)"}
}
\only<3>{
\texttt{git add share/model.csv}\\
\texttt{git commit -m"Re-trained model with 52 more participants"}
}
\uncover<4>{
\texttt{git add src/main.*}\\
\texttt{git commit -m"Move main.cpp to src/"}
}

\end{frame}

\begin{frame}[fragile]{Log}
\begin{shcode}
$ git log
commit fa009cd7fca05b0b61170b20cf76a5f72b8843c2
Author: Severin Lemaignan <severin.lemaignan@plymouth.ac.uk>
Date:   Wed Feb 10 16:48:22 2016 +0000

    Move main.cpp to src/

commit aff81119459d9193c09effef1c150c4f7eac08dc
Author: Severin Lemaignan <severin.lemaignan@plymouth.ac.uk>
Date:   Wed Feb 10 16:48:02 2016 +0000

    Re-trained model with 52 more participants

commit 4113b9b6e6bbc8de532ad90153e0059cb5819de7
Author: Severin Lemaignan <severin.lemaignan@plymouth.ac.uk>
Date:   Wed Feb 10 16:47:46 2016 +0000

    Fix computation of position (float->double)
\end{shcode}


\end{frame}

\begin{frame}{}

\centering

\begin{tikzpicture}
        \gitDAG[grow right sep = 2em]{
            4113b9 -- aff811 -- fa009c;
        };

        \gitbranch
        {master}     % node name and text 
        {above=of fa009c} % node placement
        {fa009c}          % target

        \gitHEAD
        {above=of master} % node placement
        {master}          % target

\end{tikzpicture}

\end{frame}

\begin{frame}{The staging area}
    \centering
    \only<1>{
    But why do we have to manually tell Git what files to add or remove?
    }
    \only<2>{
    To help preparing nice atomic commits!\\ $\Rightarrow$ No need to commit all current changes together.
    }
    \only<3->{
        Preparing a commit consists in filling the {\Medium staging area} (or
        {\Medium index}) with the list of changes:

\begin{tikzpicture}[
    >=latex,
    every edge/.style={draw,thick,hriSec3},
    file/.style={align=left, inner sep=1pt,font=\scriptsize\tt}
]

\node at (-3.5, 1.5)[file,red] (file1) {\sout{main.cpp}};
\node [file,red,below=0.0 of file1.south east,anchor=north east] (file2) {src/main.cpp};
\node [file,temporal=<3>{red}{black}{white},below=0.0 of file2.south east, anchor=north east] (file3) {src/position.cpp};
\node [file,temporal=<3>{red}{black}{white},below=0.0 of file3.south east,anchor=north east] (file4) {src/position.hpp};
\node [file,red,below=0.0 of file4.south east,anchor=north east] (file5) {share/model.csv};


\node[gitSA,] (stagingarea) at (-2.5,3) {staging area};

    \only<3>{
        \gitDAG[grow right sep = 2em]{
            A -- B;
        };

        \gitbranch
        {master}     % node name and text 
        {above=of B} % node placement
        {B}          % target

        \gitHEAD
        {above=of master} % node placement
        {master}          % target



        \path[draw] (file3.east) edge [->, bend left,hriSec1Comp] (stagingarea);
        \path[draw] (file4.east) edge [->, bend left,hriSec1Comp] (stagingarea);

    }
    \only<4>{
        \gitDAG[grow right sep = 2em]{
            A -- B -- C;
        };

        \gitbranch
        {master}     % node name and text 
        {above=of C} % node placement
        {C}          % target

        \gitHEAD
        {above=of master} % node placement
        {master}          % target

        \node[highlighted commit] at (C) {\phantom{C}};
        \draw[resetarrows] ([xshift=-1em,yshift=-1em]stagingarea.east) to[bend
        left] (C.north west);

    }


\end{tikzpicture}


    }
\end{frame}

\begin{frame}[fragile]{To summarize...}

The first time...
\begin{shcode}
$ mkdir my_repo && cd my_repo
$ git init
\end{shcode}
Then...
\begin{shcode}
# make some changes...
$ git add <files>
$ git commit -m"<commit message>"
# make some changes...
$ git add <files>
$ git commit -m"<other commit message>"
# That's it!
\end{shcode}


\end{frame}

\begin{frame}{Branches}

    \centering

    \begin{tikzpicture}
        \onslide<1>{
            \gitDAG[grow right sep = 2em]{
                A -- B -- C;
            };

            % Branch
            \gitbranch
            {master}     % node name and text 
            {above=of C} % node placement
            {C}          % target

            % HEAD reference
            \gitHEAD
            {above=of master} % node placement
            {master}          % target

        }

        \onslide<2>{
            \gitDAG[grow right sep = 2em]{
                A -- B -- C;
            };

            % Branch
            \gitbranch
            {master}     % node name and text 
            {above=of C} % node placement
            {C}          % target

            \gitbranch
            {new-gui}     % node name and text 
            {above=of master} % node placement
            {master}          % target

            % HEAD reference
            \gitHEAD
            {above=of new-gui} % node placement
            {new-gui}          % target
        }

        \onslide<3>{
            \gitDAG[grow right sep = 2em]{
                A -- B -- C -- D -- E;
            };

            % Branch
            \gitbranch
            {master}     % node name and text 
            {above=of C} % node placement
            {C}          % target

            \gitbranch
            {new-gui}     % node name and text 
            {above=of E} % node placement
            {E}          % target

            % HEAD reference
            \gitHEAD
            {above=of new-gui} % node placement
            {new-gui}          % target
        }

        \onslide<4>{
            \gitDAG[grow right sep = 2em]{
                A -- B -- C -- D -- E;
            };

            % Branch
            \gitbranch
            {master}     % node name and text 
            {above=of C} % node placement
            {C}          % target

            \gitbranch
            {new-gui}     % node name and text 
            {above=of E} % node placement
            {E}          % target

            % HEAD reference
            \gitHEAD
            {above=of master} % node placement
            {master}          % target
        }

        \onslide<5->{
            \gitDAG[grow right sep = 2em]{
                A -- B -- C -- {
                    D -- E,
                    F -- G -- H,
                }
            };

            % Branch
            \gitbranch
            {master}     % node name and text 
            {above=of H} % node placement
            {H}          % target

            \gitbranch
            {new-gui}     % node name and text 
            {above=of E} % node placement
            {E}          % target

            % HEAD reference
            \gitHEAD
            {above=of master} % node placement
            {master}          % target
        }


    \end{tikzpicture}

    \vspace{3em}
    \centering
    \only<1>{
        What if...?\\
    }
    \only<2>{
        \Large
        \texttt{git checkout -b new-gui}\\
    }
    \only<5>{
        The branch name is an alias for the tip of the current branch\\
    }
    \only<6>{
        $\Rightarrow$ branches are very cheap\\ 
        +10 of them at a given time it not uncommon\\
    }
    \uncover<4>{
        Let go back to serious stuff!\\
        \Large
        \texttt{git checkout master}
    }


\end{frame}


\begin{frame}{Merging branches}

    \centering

\begin{tikzpicture}[
    >=latex,
    every edge/.style={draw,thick,hriSec3}
]

        \onslide<1>{
            \gitDAG[grow right sep = 2em]{
                B -- C -- {
                    D -- E,
                    F -- G -- H,
                };
            };

            % Branch
            \gitbranch
            {master}     % node name and text 
            {above=of H} % node placement
            {H}          % target

            \gitbranch
            {new-gui}     % node name and text 
            {above=of E} % node placement
            {E}          % target

            % HEAD reference
            \gitHEAD
            {above=of master} % node placement
            {master}          % target

            \node[DAGcommit,right=3 of E.south,dashed] (merge) {?};
            \path[draw] (E) edge [->, bend left] (merge);
            \path[draw] (H) edge [->, bend right] (merge);
    }
        \onslide<2>{
            \gitDAG[grow right sep = 2em]{
                B -- C -- {
                    D -- E,
                    F -- G -- H -- merge,
                };

                E -- merge;
            };

            % Branch
            \gitbranch
            {master}     % node name and text 
            {above=of merge} % node placement
            {merge}          % target

            \gitbranch
            {new-gui}     % node name and text 
            {above=of E} % node placement
            {E}          % target

            % HEAD reference
            \gitHEAD
            {above=of master} % node placement
            {master}          % target

    }


    \end{tikzpicture}

    \vspace{3em}
    \centering
    \only<1>{
        Two options: {\Medium merging} and {\Medium rebasing}
    }
    \uncover<2>{
        Merging\\
        \Large
        \texttt{git merge new-gui}
    }


\end{frame}


\section{Sharing}

\imageframe{dvcs-1}
\imageframe{dvcs-2}
\imageframe{dvcs-3}
\imageframe{dvcs-4}
\imageframe{dvcs-5}
\imageframe{dvcs-6}
\imageframe{dvcs-7}
\imageframe{dvcs-8}


\section{GIT}

\begin{frame}[fragile]{General workflow}

The first time...
\begin{shcode}
$ mkdir my_repo && cd my_repo
$ git init
\end{shcode}
Then...
\begin{shcode}
# make some changes...
$ git add <files>
$ git commit -m"<commit message>"
# make some changes...
$ git add <files>
$ git commit -m"<other commit message>"
# ...
# when you are happy:
$ git pull # any changes on the remote?
$ git push
\end{shcode}


\end{frame}

\begin{frame}[fragile]{}

\centering

    \begin{tikzpicture}
    % Commit DAG
    \onslide<1-3>{
      \gitDAG[grow right sep = 2em]{
        A -- B -- C;
      };

      % Branch
      \gitbranch
        {master}     % node name and text 
        {above=of C} % node placement
        {C}          % target
    }



      % HEAD reference
    \onslide<1-2>{
      \gitHEAD
        {above=of master} % node placement
        {master}          % target
    }

    \onslide<3>{
      \gitremotebranch
        [origmaster]    % node name
        {origin/master} % node text
        {above=of master}    % node placement
        {master}             % target
      \gitHEAD
        {above=of origmaster} % node placement
        {origmaster}          % target
    }

    \end{tikzpicture}

\vspace{3em}
\centering
\only<2>{
\small
\texttt{git remote add origin git@github.com:user/repo.git}
}

\uncover<3>{
\Large
\texttt{git push origin master} \par
\normalsize
(or simply \texttt{git push})
}

\end{frame}

\begin{frame}[fragile]{}

    \centering

    \begin{tikzpicture}
        \onslide<1-2>{
            \gitDAG[grow right sep = 2em]{
                A -- B -- C -- D -- E;
            };

            % Branch
            \gitbranch
            {master}     % node name and text 
            {above=of E} % node placement
            {E}          % target

            % HEAD reference
            \gitHEAD
            {above=of master} % node placement
            {master}          % target

            \gitremotebranch
            [origmaster]    % node name
            {origin/master} % node text
            {above=of C}    % node placement
            {C}             % target
        }

        \onslide<3>{
            \gitDAG[grow right sep = 2em]{
                A -- B -- C -- {
                    F,
                    D -- E,
                }
            };

            % Branch
            \gitbranch
            {master}     % node name and text 
            {above=of E} % node placement
            {E}          % target

            % HEAD reference
            \gitHEAD
            {above=of master} % node placement
            {master}          % target

            \gitremotebranch
            [origmaster]    % node name
            {origin/master} % node text
            {above=of F}    % node placement
            {F}             % target
        }

        \onslide<4->{
            \gitDAG[grow right sep = 2em]{
                A -- B -- C -- {
                    F -- D' -- E',
                    {[nodes=unreachable] D -- E },
                }
            };

            % Branch
            \gitbranch
            {master}     % node name and text 
            {above=of E'} % node placement
            {E'}          % target

            % HEAD reference
            \gitHEAD
            {above=of master} % node placement
            {master}          % target

            \gitremotebranch
            [origmaster]    % node name
            {origin/master} % node text
            {above=of F}    % node placement
            {F}             % target
        }


    \end{tikzpicture}

    \vspace{3em}
    \centering
    \only<2>{
        What happened on our remote? Let's have a look...\\
        \Large
        \texttt{git fetch origin}
    }

    \only<4>{
        \Large
        \texttt{git rebase origin/master} \par
        \normalsize
        (but you don't need it, because...)
    }
    \uncover<5>{
        \Large
        \texttt{git pull --rebase} \par
    }

\end{frame}


\begin{frame}{}

    \centering

    \only<1>{
        \begin{tikzpicture}
            % Commit DAG
            \gitDAG[grow right sep = 2em]{
                A -- B -- { 
                    C,
                    D -- E,
                }
            };
            % Tag reference
            \gittag
            [v0p1]       % node name
            {v0.1}       % node text
            {above=of A} % node placement
            {A}          % target
            % Remote branch
            \gitremotebranch
            [origmaster]    % node name
            {origin/master} % node text
            {above=of C}    % node placement
            {C}             % target
            % Branch
            \gitbranch
            {master}     % node name and text 
            {above=of E} % node placement
            {E}          % target
            % HEAD reference
            \gitHEAD
            {above=of master} % node placement
            {master}          % target
        \end{tikzpicture}
    }

    \only<2>{
        \begin{tikzpicture}
            \gitDAG[grow right sep = 2em]{
                A -- B -- { 
                    C -- D' -- E',
                    {[nodes=unreachable] D -- E },
                }
            };
            % Tag reference
            \gittag
            [v0p1]       % node name
            {v0.1}       % node text
            {above=of A} % node placement
            {A}          % target
            % Remote branch
            \gitremotebranch
            [origmaster]    % node name
            {origin/master} % node text
            {above=of C}    % node placement
            {C}             % target
            % Branch
            \gitbranch
            {master}      % node name and text 
            {above=of E'} % node placement
            {E'}          % target
            % HEAD reference
            \gitHEAD
            {above=of master} % node placement
            {master}          % target
        \end{tikzpicture}
    }
\end{frame}



\begin{frame}{Working with branches}
\end{frame}

\section{Social coding: GitHub workflow}

\imageframe{../figs/github-workflow.pdf}
\imageframe{../figs/github-workflow-2.pdf}
\imageframe{../figs/github-workflow-3.pdf}

\section{Social coding: a few unwritten rules}

\begin{frame}{}
    \centering

    {\Medium Principle of least surprise}

    Make people feel at home (and save time!)

\end{frame}
\begin{frame}{Source code organization}



\end{frame}

\begin{frame}{Build system}

    \begin{itemize}
        \item P
    \end{itemize}
    
    In case of doubt, use
\end{frame}

\begin{frame}{Other tips}
    \begin{itemize}
        \item Have a \texttt{README} (or better, \texttt{README.md}): what is
            it about? who is the target audience? how to install? how to get started?
        \item No binaries!
        \item First line of commit messages < 72 characters
    \end{itemize}
\end{frame}

\begin{frame}{Software management 101}

\end{frame}

\end{document}






